\documentclass[11pt]{article}
\usepackage[margin=1in]{geometry}
\usepackage{booktabs}
\usepackage{longtable}
\usepackage{amsmath}
\usepackage{hyperref}
\hypersetup{colorlinks=true, linkcolor=blue, urlcolor=blue}

\title{SF Planning Commission Minutes:\\ Pipeline Status and Preliminary Results}
\author{Dan Post}
\date{August 2026}

\begin{document}
\maketitle

\section{Purpose}

This directory documents the empirical backbone of the project: turning twenty-eight years
of San Francisco Planning Commission meeting minutes (1998--2026) into an \emph{item-level}
dataset of discretionary land-use decisions. One row is one request heard by the
Commission --- a conditional use, a discretionary review, a variance, a rezoning --- with
its case number, address, zoning district, staff recommendation, disposition, roll-call
vote, and hearing date.

That dataset is what makes the ``market for housing regulation'' measurable. It lets us
ask who obtains discretionary approval, how often staff recommendations are overridden,
how the approval rate moves with the political composition of the Commission, and how any
of this responds to state preemption. The theory side lives in
\texttt{output/political\_economic\_housing\_model/}; the extension of the same pipeline
to the rest of the Bay Area is documented in the recon directories.

\section{What is in this directory}

\begin{longtable}{p{0.30\textwidth}p{0.62\textwidth}}
\toprule
File & What it is \\
\midrule
\endhead
\texttt{minutes\_data\_availability.md} & Inventory of what the raw files actually contain,
by document era (1{,}099 source files). Produced 2026-06-05; read-only. \\
\texttt{processing\_review.md} & Code review of the pipeline plus an audit of the
hand-labels, with a 2026-06-08 status update. \\
\texttt{data\_infrastructure.md} & The schema and how the pieces fit together. \\
\texttt{labeling\_rules.md} & The authoritative coding manual --- the per-field spec a
labeler works from. \\
\texttt{hand\_label\_review\_guide.md} & How to review existing labels before they become
ground truth. \\
\texttt{schema\_enrichment\_recommendation.md} & Field-by-field verdicts from the schema
enrichment investigation (adopt / defer / reject). \\
\texttt{meeting\_level\_info.tex} & How each item is matched to the meeting it was heard at,
the meeting-level fields, their accuracy against the gold set, and summary statistics. \\
\bottomrule
\end{longtable}

\section{The pipeline}

\textbf{Scrape} $\to$ \textbf{parse} $\to$ \textbf{date} $\to$ \textbf{label} $\to$
\textbf{extract}. Code lives in \texttt{code/commission\_minutes\_processing/}; the corpus
lives outside the repo on Dropbox and is reached through \texttt{paths.py}.

\emph{Parsing} splits each source document into blocks, one per agenda item, delimited by
case headers, numbered agenda items, section headers, adjournment markers, and
draft-minutes-adoption appendages. Two document eras are handled: scraped HTML
(1998--2014) and PDF (2015--present).

\emph{Dating} is a separate stage, and deliberately so. A block never states its own
hearing date --- the meeting header above it does --- so the date is a structural question
about which meeting a block falls under, not a textual one about the block. It is answered
by \texttt{assign\_meeting\_dates.py} and validated against hand-marked meeting boundaries
collected in \texttt{date\_boundary\_app/}. See the companion memo,
\texttt{meeting\_level\_info.tex}.

\emph{Meeting-level labelling} sits above the items. Time of day, whether the sitting was
regular, special, joint or closed, the room, who sat and who staffed it are properties of
the hearing, shared by every item heard at it, so they are recorded once per meeting and
joined on the date rather than repeated on 23{,}000 rows.
\texttt{meeting\_headers.py} cuts the header window around each known meeting start and
pre-fills twelve such fields for confirmation in the boundary app.

\emph{Labeling} is done in a local Flask app against the shared schema, with a machine
pre-fill the human corrects. \emph{Extraction} then either fine-tunes a model on those
labels or few-shot prompts one, and \texttt{learning\_curve.py} plots field accuracy
against label count so labeling stops at the empirical plateau rather than a guess.

\subsection{The schema is the single source of truth}

\texttt{extraction\_common.py} holds one list, \texttt{SCHEMA}, from which the labeling
form, the model prompt, the training builder's required keys, and the evaluation metric
are all generated. Editing that list propagates everywhere. It currently carries
\textbf{28 fields} in six groups: Identity, Location, Zoning \& scale, Process, Politics,
Conditions.

Four fields have been dropped as the schema settled, each for a stated reason:
\texttt{jurisdiction} (constant), \texttt{supervisorial\_district} (recoverable from the
address), \texttt{item} (agenda number, no analytic signal), and --- as of 30 August 2026
--- \texttt{meeting\_date}, which is a property of the meeting rather than the item and is
now supplied by the dating stage.

Two design decisions in the schema are worth stating because they carry the empirics.
First, \texttt{action} records the disposition \emph{family} (\texttt{approved},
\texttt{disapproved}, \texttt{continued}, \texttt{took\_dr\_and\_approved}, \ldots), with
``with conditions'' and ``as modified'' split out into two orthogonal flags. This removed a
three-way overlap in the old vocabulary and lets conditions attach to non-approval
outcomes. Second, the staff recommendation is stored twice: verbatim
(\texttt{preliminary\_recommendation}) and bucketed to the same disposition families
(\texttt{preliminary\_recommendation\_category}). The gap between that category and
\texttt{action} \emph{is} the staff-override signal the project studies, and storing both
makes ``was the staff recommendation followed?'' a direct enum comparison rather than
string wrangling.

\section{Preliminary results}

\subsection{Corpus}

\begin{table}[h]
\centering
\begin{tabular}{lr}
\toprule
Quantity & Count \\
\midrule
Parsed items, 1998--2026 & 23{,}043 \\
\quad HTML era (1998--2014) & 16{,}024 across 724 pages \\
\quad Modern era (2015--2026) & 7{,}019 \\
Meetings identified, HTML era & 818 \\
Pages holding more than one meeting & 41 (max 5) \\
\bottomrule
\end{tabular}
\end{table}

\subsection{Dating}

The dating stage corrected \textbf{1{,}937 of 16{,}024} HTML-era item dates. Distinct
hearing dates per year moved from 12--17 to 35--47 for 1998--2000, in line with the
Commission's roughly weekly sitting; later years were already close to correct. Against a
hand-marked gold standard of 97 documents and 101 meeting boundaries, boundary precision
and recall are both 1.000 and item date accuracy is 1.000 on 1{,}322 scored items --- and
on the 21 documents marked after the rules were frozen it is 1.000 on 334 items, so the
result is genuinely out-of-sample. All 280 modern PDFs agree with their file names. Details
and caveats are in the companion memo.

\subsection{Meetings}

All \textbf{102 meetings} of the gold sample are hand-confirmed --- the 101 marked
boundaries plus one the dating stage found that the marking missed. Against those labels
the ten meeting-level fields agree on \textbf{94.6\%} of values, and on the 21 meetings
labelled after the rules were frozen they agreed on \textbf{92.1\%} before correction, so
the extraction generalises. The sample is already informative about the panel's
composition: 90 regular sittings, 7 special, and 5 joint --- with the Redevelopment Agency
Commission (1998), Building Inspection (2007, 2018), Public Health (2011) and Historic
Preservation (2023). Two joint sittings are headed simply ``Special Meeting'' and are
identifiable as joint only from the two bodies named above the date.

\subsection{Labeling}

\begin{table}[h]
\centering
\begin{tabular}{lrl}
\toprule
Status & Items & Meaning \\
\midrule
\texttt{done} & 119 & confirmed gold \\
\texttt{review} & 80 & existing labels with a specific QA issue to recheck \\
\texttt{flagged} & 179 & 145 representativeness-sample draws + 34 recovered labels \\
\texttt{todo} & 22{,}665 & unlabeled \\
\bottomrule
\end{tabular}
\end{table}

The queue is ordered rare-class-first and interleaved by year, so labeling top-down fills
scarce request types and dispositions evenly across the panel rather than exhausting 1998
first. A fixed random sample per year is flagged separately so the gold set is
representative of the whole period rather than of whichever years were labeled earliest.

\section{Status and open questions}

\begin{itemize}
\item \textbf{Settled.} The schema; the parser's block boundaries; the dating stage and its
validation; the labeling workflow and coding manual.
\item \textbf{In progress.} Item-level hand-labeling. The immediate queue is the 80
\texttt{review} items, then the 145 representativeness draws. The meeting-level gold set
(97 documents, 102 meetings) is complete.
\item \textbf{Open.} Forty-two labels carry a hand-typed date that disagreed with the
dating pipeline (preserved in \texttt{date\_field\_disagreements.csv}); several look like
labels attached to the wrong block by the earlier content-matching recovery and need a
separate pass. The learning curve has not yet been run at full size, so the number of
labels required for the extraction model is still an estimate rather than a measurement.
\item \textbf{Decision pending.} Whether to extend meeting-level labelling beyond the
81-meeting pilot to all 818 detected meetings. The pre-fill is good enough that the cost is
confirmation rather than transcription, but the analytic payoff of the meeting-level fields
is not yet established.
\item \textbf{Known non-transfer.} The heuristic extractor is tuned to SF minutes and does
not transfer to other jurisdictions; extending the pipeline requires re-testing extraction
on real documents with human review, not unattended.
\end{itemize}

\end{document}
